% 第 49 章骨架：导引框已定，正文待写。写作计划见 docs/PLAN.md“第 49 章”，写法见 docs/STYLE.md。
\chapter{规则、可信度与沃尔克反通胀}\label{ch:credibility}

\begin{chapterguide}
\guidequestion{为什么要把印钞权交给不经选举的人？}
\begin{guidegoals}
\item 推导基德兰德与普雷斯科特（1977）、巴罗与戈登（1983）的时间不一致模型；
\item 解释罗戈夫（1985）“保守的央行行长”的含义；
\item 用沃尔克反通胀（1979—1982）估算牺牲率，并讨论可信度的作用；
\item 概括央行独立的制度设计与争论。
\end{guidegoals}
\guideprereq{\chref{ch:ratex}\enspace\nameref{ch:ratex}；\chref{ch:monetarism}\enspace\nameref{ch:monetarism}}
\end{chapterguide}

\chapterstub
